\paragraph{Measurement on shared hardware.}
All experiments ran on shared GPU nodes. Identical code drifted by up to $49\%$ between sessions, so we
build every arm once, time them round-robin in a single session, wait for an idle GPU, log other users'
GPU processes at the start and end of every repetition, and report medians over clean repetitions only.
Two complete timing passes were discarded because other users' jobs occupied the GPU. Even so,
detection is imperfect: in one Cholla pass three repetitions drifted by $15\%$ while our detector
reported no foreign GPU process, so we additionally require repetitions to agree before using them.
Residual interference cannot be excluded on a shared machine, and small effects
($\lesssim 2\%$) should not be read as real.

\paragraph{Sampling of model outputs.}
Except where stated, each (application, model, prompt) cell is a \emph{single} LLM sample. Language
model outputs vary between calls, so a single cell conflates prompt effect with sampling variance. We
replicated two settings---a second independent generation round for the structural/diagnostic arms, and
a second profile-data sample for the real applications---and the direction of the effect was stable
there, but the per-cell numbers in Section~\ref{sec:results} carry no confidence interval. Effects
below roughly $5\%$ should be treated as unresolved.

\paragraph{Hardware and software scope.}
Results cover NVIDIA GPUs only: A100 (sm\_80) for the prompt ablation, GH200 and H100 PCIe (sm\_90) for
the real applications, with CUDA 12.6.3. Conclusions may not transfer across architectures: a
diagnostics-prompt gain of $1.37\times$ on \texttt{bincount} measured on GH200 did not reproduce on the
H100 ($0.99\times$). We use two proprietary models whose weights and versions we do not control; the
same prompt may behave differently after a model update.

\paragraph{Correctness gates are empirical, not proofs.}
Each application is gated by its own output: bit-exact field comparison where the application is
deterministic (LAMMPS, Cholla) and a statistical band where it is not (AutoDock-GPU, whose GPU results
vary run to run even with a fixed seed). We validated every gate by injecting faults, and the Cholla
gate detects a $10^{-6}$ relative perturbation of a single flux. Nevertheless a subtle accuracy loss
that stays inside the AutoDock band would pass, gates check the outputs the application prints rather
than all internal state, and undefined behaviour can pass one build and fail another---our own
hand-tuned LAMMPS variant did exactly that, producing correct physics in one build and wrong physics in
a rebuild of the same source. We therefore gate every build that is timed, but ``gated correct'' means
``passed these checks'', not ``verified''.

\paragraph{Benchmark gaming is detected heuristically.}
Transformations that delete the measured work while preserving printed output defeat output comparison
by construction (Section~\ref{sec:intent}). We flag speedups above $5\times$ and inspect the diff. This
catches the large cases we observed ($1504\times$, $2809\times$, $235$--$949\times$, $31\times$), but a
gaming transformation yielding, say, $1.5\times$ would not be flagged automatically, and we cannot rule
out such cases in the reported numbers.

\paragraph{Single-file editing.}
The pipeline gives the model one source file and replaces that file. The model therefore cannot change
headers, launch sites or build flags, and---more importantly---cannot see how its file is compiled.
This caused three of our real-application failures (symbol redefinition in a textually included file,
and kernel-signature changes that broke linkage against ten call sites). Some optimizations a human
would make are out of reach by construction.

\paragraph{Tooling defects discovered during the study.}
Several defects in our own summarizers were found while running the real applications, and earlier
prompts were generated with them: shared-memory sizes were reported without their unit (a kernel using
$7.43$\,Kbyte appeared as ``7 bytes''); a PC-sampling parser dropped rows whose section header differed,
truncating the stall data for three benchmarks (for \texttt{xsbench}, 36 of 76 rows survived); true
kernel launch counts never matched namespaced or templated kernel names, so aggregate cost was
understated; comment stripping removed license headers from every generated file; and the build timeout
was reported to the model as a fault in its own code. These are fixed, but results predating the fixes
were produced with the affected prompts and would need regeneration to be strictly comparable.

\paragraph{Coverage of the profile-data ablation.}
The comparison in Section~\ref{sec:profabl} covers 26 (application, model) pairs, of which two
benchmarks (\texttt{boxfilter}, \texttt{ccsd-trpdrv}) could not be timed on the available node because
the \emph{unmodified} applications exceed our per-run cap, and two others (\texttt{adv},
\texttt{background-subtract}) fail their own built-in check in the unmodified state on that node, so
their gate carries no information there. The remaining denominator is small, and the positive result
rests on one application (Cholla) with two models.

\paragraph{Scope of the objective.}
We optimize execution time of a single GPU kernel file under a correctness gate. We do not study
energy, memory footprint, multi-GPU scaling, numerical stability under long runs, or maintainability of
the generated code---all of which matter for production use of these applications.
